\documentclass[11pt]{article}
\usepackage{fontspec}
\usepackage[french,english]{babel}
\usepackage{amsmath}
\usepackage{amssymb}
\usepackage{amsthm}
\usepackage{mathrsfs}
\usepackage{tikz-cd}
\usepackage{enumitem}

\newcommand{\printedpage}[1]{\marginpar{\footnotesize p.~#1}}
\newcommand{\msp}[1]{\marginpar{\footnotesize [#1]}}

\newcommand{\Ind}{\operatorname{Ind}}
\newcommand{\Filt}{\operatorname{Filt}}
\newcommand{\card}{\operatorname{card}}
\newcommand{\red}{\operatorname{\text{red}}}
\newcommand{\Ob}{\operatorname{Ob}}
\newcommand{\Hom}{\operatorname{Hom}}
\newcommand{\uHom}{\underline{\operatorname{Hom}}}
\newcommand{\End}{\operatorname{End}}
\newcommand{\Ens}{\mathrm{Ens}}
\newcommand{\cM}{\mathcal{M}}
\newcommand{\cD}{\mathcal{D}}
\newcommand{\defeq}{\overset{\text{def}}{=}}

\begin{document}

\printedpage{74}
% head: A. Grothendieck : Dérivateurs, Ch. XVIII, Édition des Universités de Montpellier II et Paris VII

\textbf{Lemma 4.8.10.} \textit{Let $C$ be a small category, $\pi_0$ an infinite cardinal, $\pi_1 = \sup(\pi_0, \card\red C)$. Suppose $\Ind_{\pi_0}$ is stable under finite $\varprojlim$ (resp. and of a certain type). Then for every $\pi \geqslant \pi_1$, $\Filt^{\pi}$ is stable under the said $\varprojlim$ . (Thus also $\cM_{\pi} = \Filt^{\pi}$, if $\pi = \pi^{\pi_0}$.) More} % [sic] \Ind_{\pi_0} stable -> \Ind_{\pi_0}C stable, the argument C appears to be omitted (cf. the following sentence)
\msp{page 209}
\textit{generally, let $J$ be a small category such that $\Ind_{\pi_0}C$ is stable under $\varprojlim$ of type $J$, and let $c = \card\red J$, suppose that $\pi^{c} = \pi$ (e.g. $c$ finite). Then $\Filt^{\pi}$ is stable in $\cM$ under $\varprojlim$ of type $J$. (Thus the same holds for $\cM_{\pi} = \Filt^{\pi}$, if in addition $\pi^{\pi_0} = \pi$.)}

We use criterion a$''$ of 4.7.1, namely $\card F(x) \leqslant \pi$ for every $x \in C$, for membership in $\Filt^{\pi}$. We may suppose $J$ reduced, hence $\card\Ob J \leqslant c$. Since the inclusion functor $\Ind_{\pi_0}(C) \longrightarrow C^{\wedge}$ commutes with arbitrary $\varprojlim$, the computation of $\varprojlim\limits_{j\in J} x_j$ in $C^{\wedge}$ fiber by fiber shows that in $C^{\wedge}$, we have
\[
X \ \defeq\ \varinjlim_{j\in J} x_j \ = \ \underbrace{\prod_{j\in\Ob J} x_j}_{\text{product in } C^{\wedge}} \quad ,
\] % [sic] \varinjlim_{j\in J} x_j -> \varprojlim_{j\in J} x_j, the text concerns the projective limit (cf. the preceding sentence and (4.8.10.1))
whence
\[
\card X(x) \ \leqslant\ \prod_{j\in\Ob J} \underbrace{\card x_j(x)}_{\leqslant\pi} \ \leqslant\ \pi^{c}\ ,
\]
whence (with no other hypothesis on $\pi$ than $\pi \geqslant \pi_1$)
\[
\varprojlim_{j\in J} x_j \ \in\ \Ob\Filt^{\pi^{c}}\ , \tag*{(4.8.10.1)}
\]
whence the desired conclusion if $\pi^{c} = \pi$, q.e.d.

\msp{page 210}

We are going to prove the additional implications
\[
\text{b}'' \ \Longrightarrow\ \text{b}\ , \qquad \text{c} \ \Longrightarrow\ \text{b}\ , \qquad \text{c}'' \ \overset{?}{\Longrightarrow}\ \text{c}\ ,
\]
which completes the proof.

The implication b$''$ $\Longrightarrow$ a is contained in 4.6.3.

For the implication c $\Longrightarrow$ b, we note that the stability of $\cM$ under small $\varinjlim$ is equivalent to the assertion that every small inverse limit of representable functors $\cM \xrightarrow{\ \xi_i\ } \Ens$ is representable. Now the objects of $\cM$ being accessible ($\cM$ being accessible), the representable functors $\xi_i$ are accessible, hence so is their $\varprojlim$ (SGA 4 I 9.7). On the other hand, the functor $\xi$ is left exact, as a $\varprojlim$ of representable hence left exact functors. Thus the fact that $\xi$ is representable follows from the % [sic] exactes -> exacts, agreement with « foncteurs »

\textbf{Lemma 4.8.11.} \textit{Let $\cM$ be an accessible category, stable under finite $\varprojlim$, $\cM \longrightarrow$} Ens \textit{a functor. For $\xi$ to be representable, it is necessary and sufficient that it be left exact and accessible. Thus (SGA 4 I 8.12.8.1) if $\cM$ is stable under small $\varprojlim$ , $\xi$ is representable if and only if it commutes with the said $\varprojlim$ , and is accessible.}

% foot: Version 17 décembre 2025

\printedpage{75}
% head: A. Grothendieck : Dérivateurs, Ch. XVIII, Édition des Universités de Montpellier II et Paris VII

\msp{page 211}

The ``it is necessary'' is included as a reminder, we have just seen it as a consequence of the fact that every object of $\cM$ is accessible. The ``it is sufficient'' is essentially SGA 4 I 8.13.2.b), except that in loc. cit. one supposes $\cM$ equipped with an fc that is \textit{strict.} But the proof given in loc. cit. also shows for an fcga, choosing in loc. cit. (p. 137) $\pi$ such that $\xi$ is $\pi$-accessible, and moreover $\pi$ is useful for the generalized cardinal filtration under consideration.

It remains to prove c$''$ $\overset{?}{\Longrightarrow}$ c, subject to this really being true. This would be a consequence of the following lemma (where $\cD$, $\pi$ are replaced by $C$, $\pi_0$).

\textbf{Lemma 4.8.12} (?) (\footnote[88]{Lemma not proved here -- cf. proof 4.11.7 (unfinished).}). \textit{Let $C_0$ be a small category, $\pi_0$ an infinite cardinal. If $C$ is stable under finite $\varprojlim$, then $\Ind_{\pi_0}C$ as well.} % [sic] C_0 -> C, the statement then speaks only of C

The analogous assertion for $\Ind(C)$, instead of $\Ind_{\pi}C$, is established in SGA 4 I 8.9.1 c). But it is not clear without more that $\Ind_{\pi}C$ is stable in $\Ind(C)$ under finite $\varprojlim$, in the case under consideration. More precisely, if $F = \varprojlim\limits_{j\in J} F_j$, the $F_j$ in $\Ind_{\pi}C$,
\msp{page 212}
$F$ in $\Ind(C) \subset C^{\wedge}$, we know that $C/F$ is filtering and we have to see that it is filtering large relative to $\pi_0$. We are reduced to the three cases where $J = \emptyset$, $J = e \amalg e$, and $J = \vcenter{\hbox{\begin{tikzcd}[ampersand replacement=\&, column sep=small, row sep=small]
\bullet \arrow[dr] \& \& \bullet \arrow[dl] \\
\& \bullet \&
\end{tikzcd}}}$. The first two cases are essentially trivial (for $e$, and a product of two categories large relative to $\pi_0$, are large relative to $\pi_0$), the third case on the other hand does not look trivial a priori. This is visibly related to the question of indicial representations of functors $J \longrightarrow \Ind_{\pi_0}(C)$, for $\card\red J \leqslant \pi_0$ -- i.e. the validity of 4.4.13 b) outside the hypothesis that $C$ is stable under $\varinjlim$ of cardinality $\leqslant \pi_0$. We will have to come back to this very soon, for this question now appears to me to be a key question, for the stability properties of the notion of accessible category (without any hypothesis of existence of $\varinjlim$ or $\varprojlim$ ). % [sic] qu'il est filtrante -> qu'il est filtrant, agreement with « il » (C/F)

% marginal handwritten mark (curved stroke) in the left margin beside Proposition 4.8.13
\textbf{Proposition 4.8.13} (\footnote[89]{\textit{Not proved.} But OK if $\pi_0 > \card\red C$, cf. 4.15.10, p. 462.}). \textit{Let $C$ be a category equivalent to a small category, and let $\pi_0 \geqslant \card\red C$ be an infinite cardinal. Then $\cM = \Ind_{\pi_0}C = \cM_{\pi_0}$ (i.e. all the objects of $\Ind_{\pi_0}C$ are $\pi_0$-accessible), and $\cM$ is a Karoubi envelope of $C$.}

\msp{page 213}

Let us point out at once the following corollaries.

% marginal handwritten mark (curved stroke) in the left margin beside Corollaire 4.8.14
\textbf{Corollary 4.8.14} (\footnote[90]{\textit{Not proved.} Yes after all, cf. 4.15.10, and 4.15.11, pages 462-463.}). \textit{Let $\cM$ be a $\mathfrak{U}$-category. Equivalent conditions:}
\begin{enumerate}[leftmargin=*]
\item[\textit{a)}] \textit{$\cM$ is accessible, and is equivalent to a small category.}
\item[\textit{a$'$)}] \textit{$\cM$ is accessible, and there exists an infinite cardinal $\pi_0$ such that all the objects of $\cM$ are $\pi_0$-accessible (i.e. $\cM = \cM_{\pi_0}$).}
\item[\textit{b)}] \textit{$\cM$ is equivalent to a small category, and is Karoubian, i.e. stable under direct factors} (\footnote[91]{\textbf{N.B.} This implies that $\cM$ and $\cM^{\circ}$ are accessible, and I conjecture that the converse is true as well.}). % [sic] sont accessible -> sont accessibles, plural agreement (footnote 91)
\end{enumerate}

The equivalence of a) and a$'$) is clear, since $\cM$ admits a canonical generalized cardinal filtration $(\cM_{\pi})_{\pi\geqslant\pi_0}$, consisting of categories equivalent to small categories (condition
% foot: Version 17 décembre 2025
\printedpage{76}
% head: A. Grothendieck : Dérivateurs, Ch. XVIII, Édition des Universités de Montpellier II et Paris VII
a) for fcg), and such that every such subcategory is contained in one of the $\cM_{\pi}$ (follows from condition b) for fcg). Since $\cM_{\pi_0}$ is equivalent to a small category, and is Karoubian (being stable under filtering $\varinjlim$ large relative to $\pi_0$), it is also clear that a$'$) $\Longrightarrow$ b). It remains to prove b) $\Longrightarrow$ a$'$), now this follows at once from 4.8.13 applied to $\cM$, which moreover shows that for $\pi_0$ one can take any infinite cardinal such that $\pi_0 \geqslant \card\red\cM$.

\msp{page 214}

\textbf{4.8.15.} \textsc{Proof of} 4.8.13. We are going to prove:

\textbf{Lemma 4.8.15.1.} \textit{Under the conditions of 4.8.13, let $F \in \Ob\cM$, $x \in \Ob C$, then}
\[
\card F(x) \ \leqslant\ \pi_0\ . \tag*{$(*)$}
\]

This therefore means, by 4.7.7, a$''$ $\Longrightarrow$ b for $\pi \geqslant \pi_1$, taking into account that here $\pi_J = \pi_0$, $\pi = \pi_0$, that we have $F \in \Filt^{\pi_0}$, whence a fortiori $F \in \cM_{\pi_0}$, which implies $\cM = \cM_{\pi_0}$, i.e. proposition 4.8.13, taking 4.8.5 into account. So it all comes down to proving $(*)$.

Let us write
\[
F \ = \ \varinjlim_{i\in I} x_i\ , \qquad\quad \text{the } x_i \in \Ob C,\ I \text{ ordered set large relative to } \pi_0,
\]
whence
\[
F(x) \ = \ \varinjlim_{i\in I} \Hom_C(x, x_i)\ .
\]
We may suppose
\[
\card C \ = \ \pi_0\ ,
\]
replacing $C$ by an equivalent reduced category if need be.

For every $\xi \in \Ob C$, let
\[
I_{\xi} \ = \ \{ i \in I \mid x_i = \xi \}\ ,
\]
then
\[
I \ = \ \bigcup_{\xi\in\Ob C} I_{\xi}\ ,
\]
\msp{page 215}
a union of a family of subsets of $I$ of cardinality $\card\Ob C \leqslant \pi_0$. Since $I$ is large relative to $\pi_0$, we see that this implies that the set of the $I_{\xi}$ is cofinal in $I$. Replacing $I$ by $I_{\xi}$ if need be, we may suppose $I_{\xi} = I$, i.e. that all the $x_i$ are equal to the same element $\xi$. Replacing $C$ by the full subcategory $C_0 = \{\xi\}$ if need be, we may suppose $C$ reduced to a single element. (For if we prove that $F$, as an object of $\Ind_{\pi_0}(C_0)$, is a direct factor of an object of $C_0$ (i.e. of $\xi$), the same will hold for $F$ as an object of $\Ind C$). Thus we may suppose
\[
C \ = \ \{\xi\}\ ,
\]
thus $C$ given by the monoid
\[
M \ = \ \End_C(\xi)\ , \qquad\quad \card M \ \leqslant\ \pi\ ,
\]
and giving the direct system $(x_i)_{i\in I}$ amounts to giving a system of elements
\[
\alpha_{ji} \ \in\ M \qquad\quad (i, j \in I,\ i \leqslant j)
\]
% foot: Version 17 décembre 2025
\printedpage{77}
% head: A. Grothendieck : Dérivateurs, Ch. XVIII, Édition des Universités de Montpellier II et Paris VII
satisfying
\[
\alpha_{ii} \ = \ 1\ , \qquad\quad \alpha_{kj}\alpha_{ji} \ = \ \alpha_{ki} \ \ \text{if } i \leqslant j \leqslant k.
\]
We have
\[
F(\xi) \ = \ \varinjlim_{i\in I} M_i \qquad\quad
\left\{
\begin{array}{l}
\text{where } M_i = M \ \forall\, i \in I, \text{ and } M_i \longrightarrow M_j \\
\text{is given by} \\[2ex]
\qquad\qquad u \longmapsto \alpha_{ji}u\ . \\[4ex]
\end{array}
\right.
\]

\msp{page 216}

I have been able to reduce, by similar arguments (\footnote[92]{For these reductions, cf. 4.15.8 (p. 453-463). The ultimate reduction is found in 4.15.8.9 (p. 461), and its corollary 4.15.10.}), to the case where the monoid $M = \End(\xi)$ is a monoid of monomorphisms, i.e.
\[
uv \ = \ uv' \ \ \Longrightarrow\ \ v \ = \ v' \qquad\quad \text{if } u, v, v' \in \cM\ ,
\] % [sic] u, v, v' \in \cM -> u, v, v' \in M, the elements belong to the monoid M
which implies that the transition arrows in the direct system of the $M_i = M$ are monomorphisms. It follows that one can have $\card F(\xi) \leqslant \pi_0$ ($I$ being large relative to $\pi_0$) only if $\alpha_{ji_0}$ is an isomorphism for $j \geqslant i_0$, and $i_0$ large enough, i.e. if the induced direct system for $j \geqslant i_0$ is constant, a fortiori that the direct system under consideration is essentially constant. But I have just spent nearly four hours trying to prove this point, without succeeding! This seems more delicate to prove than I had expected, but I still believe that it really is true.

\msp{page 217}

\subsection*{4.9\quad The indicial representation theorem for $\Ind_{\pi}(C)$}

\textbf{4.9.1} (\footnote[93]{\texttt{[The numbers 4.9.1 to 4.9.6 are used twice, 4.9.3 three times.]}}). Let $C$ be a category equivalent to a small category, $\pi_0$ an infinite cardinal,
\[
\cM \ = \ \Ind_{\pi_0}(C) \ \subset\ C^{\wedge}\ , \tag*{(4.9.1)}
\]
$J$ a category equivalent to a small category. Consider the fully faithful inclusion $C \hookrightarrow \cM$, and the corresponding inclusion of $\mathfrak{U}$-categories
\[
\uHom(J, C) \ \hookrightarrow\ \uHom(J, \cM) \qquad\quad \text{(fully faithful).} \tag*{(4.9.2)}
\]
We know, since $C \subset \cM_{\pi_0}$, that the objects of $\uHom(J, C)$ give $\pi_0$-\textit{accessible} objects of $\uHom(J, \cM)$, at least if we suppose
\[
\card\red J \ \leqslant\ \pi_0\ , \tag*{(4.9.3)}
\]
which we shall do henceforth. (Cf. 4.5.8, 2°, p. 119.) As the right-hand side of (4.9.2) satisfies $\mathrm{L}_{\pi_0}$ ($\cM$ satisfies it), it follows that this functor (4.9.2) extends to
\[
\Ind_{\pi_0}(\uHom(J, C)) \ \hookrightarrow\ \uHom(J,\ \underbrace{\cM}_{=\,\Ind_{\pi_0}(C)}\ )\ , \tag*{(4.9.3)}
\]
% foot: Version 17 décembre 2025
\printedpage{78}
% head: A. Grothendieck : Dérivateurs, Ch. XVIII, Édition des Universités de Montpellier II et Paris VII
and this functor is \textit{fully faithful} thanks to the accessibility result we have just stated. I am interested in establishing cases where this functor is essentially surjective, i.e. is an \textit{equivalence} of categories. Let us note that for
\[
\varphi \ :\ J \longrightarrow \cM
\]
to belong to the essential image, it is necessary and sufficient that there exist a functor
\[
\Phi \ :\ I \ \longrightarrow\ \uHom(J, C)\ ,\ I \text{ large relative to } \pi_0 \text{ (and we may suppose } I \text{ ordered)}\ , \tag*{(4.9.4)}
\]

\msp{page 218}

such that we have
\[
\varphi \ \xleftarrow{\ \sim\ }\ \varinjlim_{i\in I} \Phi(i)\ , \tag*{(4.9.5)}
\]
or what amounts to the same thing, such that we have
\[
\varphi(j) \ \xleftarrow{\ \sim\ }\ \varinjlim_{i\in I} \Phi(i, j)\ , \qquad\quad \text{for } j \in J\ , \tag*{(4.9.6)}
\]
where we have put
\[
\Phi(i, j) \ = \ \Phi(i)(j)\ ,
\]
thus interpreting $\Phi$ as a functor
\[
\Phi \ :\ I \times J \ \longrightarrow\ C\ ; \tag*{(4.9.7)}
\]
the functorial law of $\varphi(j)$ given by (4.9.6) being evident. In other words, the $\varphi$ in the essential image of (4.4.3) are those which admit an \textit{indicial representation} by a functor $\Phi$ in the form (4.9.4), or in the form (4.9.7), with \textit{in addition} $I$ large relative to $\pi_0$. I recall in this connection the % [sic] a ce sujet -> à ce sujet, missing accent

\textbf{Proposition 4.9.2.} \textit{Under the preceding conditions on $C$, $\pi_0$, $J$ (in particular} $\card\red J \leqslant \pi_0$\textit{), if $\varphi : J \longrightarrow \cM \defeq \Ind_{\pi_0}(C)$, for $\varphi$ to belong to the essential image of the fully faithful functor (4.9.3), it is necessary and sufficient that it satisfy the following two conditions:}

\msp{page 219}

\begin{enumerate}[leftmargin=*]
\item[\textit{a)}] $\uHom(J, C)/\varphi$ \textit{is large relative to $\pi_0$.}
\item[\textit{b)}] \textit{For every $j \in J$, the functor}
\[
\uHom(J, C)/\varphi \ \longrightarrow\ C/\varphi(j) \qquad\quad
\left\{
\parbox{0.35\textwidth}{\textit{functor between filtering categories large relative to $\pi_0$, by virtue of a),}}
\right. \tag*{(4.9.2.1)}
\]
\textit{induced by the functor}
\[
\left\{
\begin{array}{rcl}
\uHom(J, C) & \longrightarrow & C \\
F & \longmapsto & F(j)\ ,
\end{array}
\right. \tag*{(4.9.2.2)}
\]
\textit{is} cofinal.
\end{enumerate}
% foot: Version 17 décembre 2025

\end{document}
